\begin{figure}[t]
\centering
\begin{tikzpicture}[
  box/.style={draw=linkblue,rounded corners=2pt,align=center,
    text width=55mm,minimum height=13mm,inner sep=5pt,font=\small},
  arr/.style={-{Stealth[length=2mm]},draw=linkblue,thick},
  node distance=10mm and 8mm]
\node[box] (minus) {First level $\ell_1$\\
  center and target planes span\\a degenerate four-space};
\node[box,right=of minus] (plus) {Second level $\ell_2$\\
  center and target planes span\\a nondegenerate four-space};
\node[box,below=of minus] (sign) {Rules out the negative sign\\
  in a Frobenius identity};
\node[box,below=of plus] (torus) {Rules out a common\\
  toric two-space at a bad center};
\draw[arr] (minus)--(sign);
\draw[arr] (plus)--(torus);
\node[box,text width=70mm,below=14mm of $(sign.south)!0.5!(torus.south)$]
  (height) {Common equations and nonidentity\\
  + proximity and coordinate bounds\\
  $\Longrightarrow$ polynomial height bound};
\draw[arr] (sign.south)--(height.north west);
\draw[arr] (torus.south)--(height.north east);
\end{tikzpicture}
\caption{The two level incidences have different roles. The first is used
to exclude the negative sign in the nonidentity argument; the positive
sign is excluded by the Frobenius multiplier. The second excludes
proximity to centers with non-potentially-good reduction. The height
conclusion uses common equations at the retained places, nonidentity on
algebraic boundary branches, the large proximity sum, and the small
positive coordinate logarithms described in the text. The diagram is
schematic: all plane positions are in the separately marked finite-field
quadratic spaces.}
\label{fig:two-incidences}
\end{figure}
